% ============================================================================
% RELATED: MANIFOLD CAPACITY AND REPRESENTATION GEOMETRY
% ============================================================================
\subsection{Manifold Capacity and Representation Geometry}
\label{sec:related_manifold_capacity}

\paragraph{Manifold Capacity Theory.}
Chung, Lee, and Sompolinsky \citep{chung2018classification} developed a
statistical-physics theory of \emph{perceptual manifold capacity}: the
maximum number of object manifolds that a linear readout can separate.
Their framework establishes that capacity depends not on ambient
dimensionality alone but on the \emph{geometry} of the manifolds---their
intrinsic dimension, radius, and center correlations.  Recent extensions
\citep{cohen2024manifoldcapacity} apply this theory across biological
neural recordings and artificial networks, introducing Geometric measures
from Correlated Manifold Capacity (GCMC) that analytically connect
neural co-variabilities to downstream task efficiency.

The Value representation provides a concrete setting in which these
geometric quantities can be characterized.  Each Value occupies a point
in $\{0,1\}^{8192}$ and also carries an eight-lane PGA coordinate in its
Context region.  A collection of Values derived from related inputs (for
example, different byte sequences sharing a common prefix) therefore
defines both a discrete affinity manifold and a continuous local
transition manifold.  The relevant geometric properties are:

\begin{itemize}[nosep]
  \item \textbf{Intrinsic dimension:} bounded by the number of bits
        that vary across the manifold, which depends on the length and
        diversity of the input sequences.
  \item \textbf{Radius:} controlled by the Hamming distance between the
        most dissimilar Values in the manifold.
  \item \textbf{Center correlations:} determined by the shared
        bit-pattern structure among Values derived from related inputs.
\end{itemize}

A population-count-based similarity score between two Values---the
Hamming inner product---functions as a linear readout over the binary
manifold space, while cosine similarity between multivectors functions
as the local readout for the geometric lane.  The discriminative
capacity of these readouts is therefore governed by the geometric
properties that Chung's theory characterizes.  Formal application of
manifold capacity bounds to the Value representation is a direction for
future analytical work.

\paragraph{Relationship to Noncommutative Algebra.}
Hataya and Hashimoto \citep{hataya2023noncommutative} provide a formal
proof that noncommutative product structures in parameter spaces induce
improved feature learning for structured tasks.  Their framework
generalizes neural network parameters from scalars to elements of a
noncommutative $C^*$-algebra.

As discussed in \Cref{sec:related_algebraic_topology}, the composition
of Boolean and PGA operations over Values is non-commutative: different
orderings of the same operations produce different frame states.  This
non-commutativity is structural rather than learned---it arises from the
sequential application of truth tables and geometric products to a shared
frame substrate.  The question of whether this prescribed
non-commutative structure yields the representational benefits that
Hataya and Hashimoto's theory predicts is addressed empirically in the
experimental sections.
